%2025 - 8a


{\bf Chapitre 5 : Dérivée covariante, connexion (suite)}


{\bf 8a Rappel}
\subsection{Dérivée covariante}

On a une variété avec une topologie et un atlas lisse(M, O, A), C-infini. A partir de cet atlas on peut définir des vitesses V et un espace vectoriel TpM (espace vectoriel tangent). Celui-ci a un dual T*pM. On construit alors l'ensemble fibré tangent TM = union disjointe des TpM dont on fait une variété différentielle. On peut alors définir un champ de vecteur comme une section du fibré tangent. On a introduit alors la dérivé covariante permettant de dériver un champ de vecteur lisse. 

Leibniz : quand plusieurs choses varient, on somme les variations.

\subsection{Exemples de connexions}
Plan euclidien à deux dimensions, R2, dans la carte identité :
 
\[
\begin{array}{ l c l l }
 id : & \mb{R}^2 & \to & \mb{R}^2 \\ 
  & (x,y) & \mapsto & (x,y) \\
\end{array} 
\]

Sur la carte identité, on peut choisir la connexion
\[
\qquad \Gamma_{(id)\ jk}^i(\lambda)=0
\]
ce qui donne la géométrie habituelle

Dans la carte polaire (r, théta)

À chaque choix correspond une géométrie particulière

On pourrait décider que les gamma de la carte polaire sont nuls, Dès lors, nous ne serions plus dans le plan euclidien.

"La dérivée covariante donne une forme à la variété".

\subsection{Transport parallèle}
Un champ de vecteur X est dit parallèlement transporté le long d'une courbe gamma I->M si la dérivée covariante dans la direction de la vitesse de la courbe en chaque point du champ de vecteur est égale à zéro.

Dans une carte (U,x )
\subsection{Courbe autoparallèle et autoparallèlement transportée}
\subsection{Équation géodésique (dans une carte donnée)}
\subsection{Annexe : reparamétrage}
